% Source: https://artofproblemsolving.com/community/c3h1456656_a_pigeonhole_problem
% Stable identity: aops-topic-1456656-q01
Take any student $X$ and let $d$ be the number of students $X$ knows ("knowing" is mutual).

\textbf{$d\ge3$ is impossible.} Let $Y_1,Y_2,Y_3$ be three acquaintances of $X$. If two of them, say $Y_1$ and $Y_2$, know each other, then in the triple $X,Y_1,Y_2$ everybody knows everybody, so it has no two students who do not know each other, a contradiction. So $Y_1,Y_2,Y_3$ pairwise do not know each other, and this triple has no two students who know each other, again a contradiction.

\textbf{$d\le1$ is impossible.} Then $X$ does not know at least $3$ students $Z_1,Z_2,Z_3$. If two of them, say $Z_1,Z_2$, do not know each other, the triple $X,Z_1,Z_2$ has no acquainted pair. So $Z_1,Z_2,Z_3$ all know each other, and that triple has no unacquainted pair, a contradiction.

Hence $d=2$ for every student. $\blacksquare$ (In fact the acquaintances then form a $5$-cycle: the students can be seated around a table so that each knows exactly the two neighbours.)
